% Compile with: uplatex file.tex && dvipdfmx file.dvi
%           or: latexmk -uplatex file.tex
\documentclass[dvipdfmx]{article}
\usepackage{amsmath, amssymb}
\usepackage{tikz}

\title{scenario\_pendulum\_lagrangian}
\date{\today}

\begin{document}
\maketitle

Consider a simple pendulum (length $l$, angle $q$, mass $m$). Prepare $q$ as a time-dependent angle, and $m,l,g$ (gravitational acceleration) as scalar parameters.

Define q as a time-dependent function

$\checkmark$ Declared q as time-dependent (time variable: t)

Define m, l, g as scalars

\[
  m, l, g
\]

Define the pendulum's Lagrangian $L=\frac12ml^2\dot q^2-mgl(1-\cos q)$ (kinetic energy minus potential energy).

$L = (1/2) m l^2 \dot{q}^2 - m g l (1 - cos(q))$

\begin{equation}
  L = (1/2) m l^2 \dot{q}^2 - m g l (1 - \cos(q))
  \label{eq:E1}
\end{equation}

Following the Euler-Lagrange equation $\frac{d}{dt}\frac{\partial L}{\partial\dot q}-\frac{\partial L}{\partial q}=0$, first find the generalized momentum $\partial L/\partial\dot q$.

Take the partial derivative of E1 with respect to $\dot{q}$

\begin{equation}
  \frac{\partial L}{\partial \dot{q}} = m l^2 \dot{q}
  \label{eq:E2}
\end{equation}

Take the total time derivative of the generalized momentum just found, to obtain $\frac{d}{dt}(\partial L/\partial\dot q)$.

Take the total derivative of E2 with respect to t

\begin{equation}
  \frac{d}{dt}\left(\frac{\partial L}{\partial \dot{q}}\right) = m l^2 \ddot{q}
  \label{eq:E3}
\end{equation}

Find the other term of the Euler-Lagrange equation, $\partial L/\partial q$.

Take the partial derivative of E1 with respect to q

\begin{equation}
  \frac{\partial L}{\partial q} = -m g l \sin q
  \label{eq:E4}
\end{equation}

Subtract $\partial L/\partial q$ from $\frac{d}{dt}(\partial L/\partial\dot q)$ to assemble the equation of motion.

$EOM = E3 - E4$

\begin{equation}
  EOM = m l^2 \ddot{q}+m g l \sin q
  \label{eq:E5}
\end{equation}

Divide both sides by $ml^2$ to drop the leading coefficient, moving toward the simple-harmonic-motion form.

Divide both sides of E5 by m $l^2$

\begin{equation}
  \frac{EOM}{m l^2} = \frac{l \ddot{q}+g \sin q}{l}
  \label{eq:E6}
\end{equation}

Apply the small-angle approximation $\sin q\approx q$, reducing the equation to simple harmonic motion $\ddot q+\omega^2q=0$ ($\omega=\sqrt{g/l}$).

In E6, substitute q for sin(q)

\begin{equation}
  \frac{EOM}{m l^2} = \frac{l \ddot{q}+g q}{l}
  \label{eq:E7}
\end{equation}

\end{document}